% Figure: effectiveness vs NTU for four flow arrangements at a fixed capacity ratio
% Cr = 0.5: counter-flow (best), cross-flow both-unmixed, one-shell-pass two-tube-pass,
% and parallel-flow (worst). The ranking is the design content: at the same area the
% arrangement sets how much of the maximum duty is reached.
\begin{figure}[htbp]
\centering
\begin{tikzpicture}
\begin{axis}[
  width=12cm, height=8cm,
  xlabel={number of transfer units, $\mathrm{NTU}=UA/C_{\min}$},
  ylabel={effectiveness $\varepsilon$ (at $C_r=0.5$)},
  xmin=0, xmax=5, ymin=0, ymax=1.0,
  axis lines=left,
  legend pos=south east,
  legend cell align=left,
  tick label style={font=\scriptsize},
  label style={font=\small},
  legend style={font=\scriptsize},
  samples=120, domain=0.001:5,
]
% Counter-flow, Cr=0.5
\addplot[engblue, very thick]
  {(1 - exp(-x*(1-0.5)))/(1 - 0.5*exp(-x*(1-0.5)))};
\addlegendentry{counter-flow}
% Cross-flow both unmixed, Cr=0.5 (closed-form correlation)
\addplot[engamber, very thick]
  {1 - exp( (exp(-0.5*x^0.78) - 1)*x^0.22/0.5 )};
\addlegendentry{cross-flow, both unmixed}
% One shell pass, 2 tube passes, Cr=0.5:
% eps1 = 2 / { (1+Cr) + sqrt(1+Cr^2) * (1+E)/(1-E) }, E = exp(-NTU*sqrt(1+Cr^2))
\addplot[engorange, very thick, samples=120, domain=0.001:5]
  ( {x}, { 2/( (1+0.5) + sqrt(1+0.25)*( (1+exp(-x*sqrt(1.25)))/(1-exp(-x*sqrt(1.25))) ) ) } );
\addlegendentry{1 shell pass, 2 tube passes}
% Parallel-flow, Cr=0.5
\addplot[engred, very thick]
  {(1 - exp(-x*(1+0.5)))/(1+0.5)};
\addlegendentry{parallel-flow}
\end{axis}
\end{tikzpicture}
\caption[Effectiveness versus NTU for four arrangements]{Effectiveness against the
number of transfer units at a fixed capacity ratio $C_r=0.5$, for four flow
arrangements that share the same energy balance and differ only in how the two
streams are routed past each other. Counter-flow is the best arrangement at every
NTU, parallel-flow the worst, with cross-flow and the one-shell-pass two-tube-pass
shell-and-tube unit between them. The spread widens as NTU grows: at small area all
four are close, because little heat is transferred and the routing scarcely
matters, while at large area the parallel-flow curve flattens against its ceiling
(parallel flow cannot drive the two outlets across each other) and the counter-flow
curve climbs on toward unity. The practical reading is that the arrangement is a
free lever at high NTU: the same tubes routed counter-current rather than co-current
buy a substantial increment of duty at no extra area, which is why counter-flow is
the default and parallel-flow is reserved for the rare service that wants its limited
approach.}
\label{fig:vol04ch06:arrangement-effectiveness}
\end{figure}
